% 第7批：《大学物理简明教程》第7章静电场解答题7.3.9—7.3.14。
% 原书/PDF第188页，模型从用户上传图核对。
\subsection{解答题：静电势、位移功和电容}

\sourceproblem{7.3.9}{188}{188}
\problemtext 均匀线电荷密度$\lambda$分布在长为$L$的细杆上，
求在杆轴延长线、距最近端$a>0$的点$P$的电势，
以无穷远处电势为零。
\workings
取杆最近端到最远端的坐标$s\in[0,L]$，
微元电荷$dq=\lambda\,ds$距$P$为$a+s$。
故
\[
\varphi_P=\int_0^L\frac1{4\pi\varepsilon_0}
\frac{\lambda\,ds}{a+s}
=\boxed{\frac{\lambda}{4\pi\varepsilon_0}
\ln\frac{a+L}{a}}.
\]
\crosscheck 当$L\ll a$时，
$\ln(1+L/a)\approx L/a$，退化为总电荷$\lambda L$
集中于远处的点电荷电势$k_e\lambda L/a$。

\sourceproblem{7.3.10}{188}{188}
\problemtext 四个相同点电荷
$q_1=q_2=q_3=q_4=1.25\times10^{-8}$ C
分别位于正方形四个顶点，每顶点到中心$O$的距离$r=5.0\times10^{-2}$ m。
(1) 求中心电势；(2) 从无穷远把试验电荷$q_0=1.0\times10^{-9}$ C搬到中心，
求静电场所做功。
\workings
电势是标量，四个电荷的贡献直接相加，
\[
\boxed{\varphi_O=\sum_{j=1}^4\frac{k_eq_j}{r}
=\frac{4k_eq}{r}
\approx8.99\times10^3\,\mathrm V}.
\]
试验电荷的电势能从零增至$q_0\varphi_O$，
静电场做功为其负值，
\[
\boxed{W_{\infty\to O}=-q_0\varphi_O
\approx-8.99\times10^{-6}\,\mathrm J}.
\]
\crosscheck 中心电场因空间对称性可以为零，
但中心电势为正且不为零；不能把“场强零”误读为“电势零”。

\sourceproblem{7.3.11}{188}{188}
\problemtext 直线上点$A,B$分别放置$+q,-q$，
$AB=2R$，中点为$O$。点$C$位于$B$右侧距$B$为$R$处。
试验电荷$q_0$从$O$经上方以$O$为起点的半圆弧移动到$C$，
求过程中的静电场力做功。
\workings
取$O$为坐标原点，
$A=-R$、$B=+R$、$C=+2R$。
在$O$点，
\[
\varphi_O=k_e\left(\frac qR-\frac qR\right)=0.
\]
在$C$点，距$A$为$3R$、距$B$为$R$，
\[
\varphi_C=k_e\left(\frac{q}{3R}-\frac qR\right)
=-\frac{2k_eq}{3R}.
\]
因此
\[
\boxed{W_{O\to C}=q_0(\varphi_O-\varphi_C)
=\frac{2k_eqq_0}{3R}
=\frac{qq_0}{6\pi\varepsilon_0R}}.
\]
\crosscheck 静电力是保守力，与所经过半圆的几何形状无关；
只需要$O,C$两端的电势。

\sourceproblem{7.3.12}{188}{188}
\problemtext 如图7.3.12，均匀正线电荷密度$\lambda$分布在同一平面内的一根弯曲绝缘细线上：
中间是上半圆、半径$R$，两端各接一段长度为$R$的共线水平直线段。
求半圆圆心$O$处的场强矢量与电势（无穷远零电势）。
\workings
设半圆圆心为原点、半圆位于$y>0$。
左右直线段分别处于$[-2R,-R]$及$[R,2R]$，
对中心的水平电场贡献大小相等、方向相反，故合场为零。
半圆弧长微元$d\ell=R\,d\theta$，
每段到圆心距离恒定为$R$。
其场强水平分量对称相消，竖直分量
\[
E_y=-\frac{k_e\lambda}{R}
\int_0^\pi\sin\theta\,d\theta
=-\frac{2k_e\lambda}{R}.
\]
所以
\[
\boxed{\bm E_O=-\frac{2k_e\lambda}{R}\bm e_y}.
\]
电势为标量积分；
半圆贡献$k_e\lambda\pi$，两水平段各贡献
$k_e\lambda\int_R^{2R}dx/x=k_e\lambda\ln2$，
故
\[
\boxed{\varphi_O=k_e\lambda(\pi+2\ln2)}.
\]
\crosscheck 仅半圆弧提供非零的总场强，
但两端直线段都提供非零\emph{电势}；
两种量叠加方式不同。

\sourceproblem{7.3.13}{188}{188}
\problemtext 一均匀带电细圆环半径为$R$、总电荷$Q$，
在环轴线上距环中心$O$为$x$的点$P$求电势。
\workings
环上任意电荷微元$dq$到$P$的距离相同，均为
$s=\sqrt{R^2+x^2}$，所以
\[
\varphi_P=\frac1{4\pi\varepsilon_0}
\int\frac{dq}{\sqrt{R^2+x^2}}
=\boxed{\frac{Q}{4\pi\varepsilon_0\sqrt{R^2+x^2}}}.
\]
\crosscheck $x=0$时$\varphi_O=k_eQ/R$；
$|x|\gg R$时退化为点电荷远场电势$k_eQ/|x|$。

\sourceproblem{7.3.14}{188}{188}
\problemtext 两个平行金属板$A,B$的正对面积为$S=200\,\mathrm{cm^2}$，
板间距$d=2\,\mathrm{cm}$。$B$板接地，$A$板带正电
$Q=7.08\times10^{-7}$ C，忽略边缘效应，取地电势为零。求$A$板电势。
\workings
面积$S=0.0200\,\mathrm{m^2}$、间距$d=0.0200$ m，
平行板电容
\[
C=\frac{\varepsilon_0 S}{d}
=\varepsilon_0
=8.85\times10^{-12}\,\mathrm F
\]
（采用题设真空介电常数的三位有效数字）。
$B$板接地$\varphi_B=0$，
\[
\boxed{\varphi_A=U_{AB}=\frac{Q}{C}
\approx8.0\times10^4\,\mathrm V}.
\]
\crosscheck 在$B$接地且$A$带正电时，
$\varphi_A>\varphi_B$；
只有忽略边缘效应的理想平行板近似下才能直接使用$C=\varepsilon_0S/d$。
